\begin{frame}
  \frametitle{Knuth-Morris-Pratt algorithm}

Knuth proposed an improvement: if sliding the pattern entails
comparing the same letter from the pattern with the same letter in the
text, \textbf{we should slide again} because the same mismatch will
happen, and again until we find a new letter.

\begin{figure}
\centering
\includegraphics{kmp_ex}
\end{figure}

\end{frame}


\begin{frame}
  \frametitle{Knuth-Morris-Pratt algorithm}

Let us see the need for this improvement in another way, by
representing graphically \(\MPfailureName_x\), with \(x =
\word{abacabac}\), as we know it:

\bigskip

\begin{figure}
\centering
\includegraphics[bb=80 651 390 710]{abacabac}
\end{figure}

\medskip

This is the Morris-Pratt \emph{deterministic finite automaton} (DFA)
for the word \word{abacabac}. The table we gave earlier was:

\begin{figure}
\centering
\includegraphics{mp_fail}
\end{figure}

\end{frame}

\begin{frame}
  \frametitle{Knuth-Morris-Pratt algorithm}

We will here only describe informally such automata.

\bigskip

Consider that the circles, called \textbf{states}, contain the values
of~\(i\) from the table. The edges, called \textbf{transitions},
between two states are of two kinds: either solid and carrying a
letter, called \textbf{label}, or dotted and going backwards. The
sequence of states throughout solid edges make the word
\(x=\word{abacabac}\).

\bigskip

The rightmost state is distinguished by a double circling because it
marks the end of~\(x\). There is a back edge between state~\(i\)
and~\(j\) only if \(\MPfailure{x}{}{i} = j\).

\bigskip

The leftmost state has an incoming, solid edge without a source and a
dotted, outgoing edge. The former simply denotes the beginning of the
word and the latter corresponds to the special value
\(\MPfailure{x}{}{0} = -1\).

\end{frame}

\begin{frame}
  \frametitle{Knuth-Morris-Pratt algorithm}

In terms of the automaton, the improvement consists in following the
back edges

\bigskip

\begin{figure}
\centering
\includegraphics[bb=82 646 388 720]{kmp_abac}
\end{figure}

\bigskip

Intuitively, it goes as follows. At a given state with the outgoing
label~\(a\), follow the back edge to another state. If the outgoing
label is~\(a\) then repeat, otherwise we are done and we process
another state. If the back edge arrives at state~\(0\) and the
outgoing label is still~\(a\), it loses its destination.

\end{frame}


\begin{frame}
  \frametitle{Knuth-Morris-Pratt algorithm}

  It terms of a table, this DFA is equivalent to

\begin{figure}
\centering
\includegraphics{kmp_fail}
\end{figure}

\begin{codebox}
\(\proc{Next'} (x)\)\\
\li \(\id{next} \gets \proc{Next}(x)\)
\li \For \(i \gets 1\) \To \(\len{x} - 1\)
\li \Do
\li \(j \gets \id{next}[i]\)
\li \If \(j \geqslant 0\) \LogAnd \(x[j] = x[i]\)
\li \Then \(\id{next}[i] \gets \id{next}[j]\) \End
\End
\li \(\proc{Next'} \gets \id{next}\)
\end{codebox}

\end{frame}

\begin{frame}
  \frametitle{Knuth-Morris-Pratt algorithm}

With some effort, we can fuse the two loops into one. The idea is to
maintain~\(j\) as the next state before optimisation.

\begin{codebox}
\(\proc{Next'} (x)\)\\
\li \(\id{next}[0] \gets -1\); \(j \gets 0\)
\li \For \(i \gets 1\) \To \(\len{x} - 1\)
\li \Do
\li   \If \(x[j] = x[i]\)
\li   \Then \(\id{next}[i] \gets \id{next}[j]\)
\li   \Else
\li     \(\id{next}[i] \gets j;\)
\li     \While \(j \geqslant 0\) \LogAnd \(x[j] \neq x[i]\)
\li     \Do
\li       \(j \gets next[j]\)
        \End
      \End
\li \(j \gets j + 1\)
      \End
\li \(\proc{Next'} \gets \id{next}\)
\end{codebox}


\end{frame}
